\section{Design Principles}
\label{sec:principles}

The following principles govern Timeline Compiler's architecture. They are listed in priority order; when principles conflict, earlier principles take precedence.

\subsection{Presentation-First}
\label{principle:presentation-first}

\textbf{Principle:} Every design decision optimizes for presentation-quality visual output.

\textbf{Rationale:} The distinguishing feature of Timeline Compiler is producing visuals suitable for board presentations, investor decks, and executive communication. This is not a documentation tool (Mermaid serves that purpose) nor a project management tool (many exist). If a feature does not improve presentation quality, it should not exist.

\textbf{Implications:}
\begin{itemize}
  \item Typography, spacing, and color receive first-class treatment in the theme engine
  \item The IR includes presentation hints (emphasis, annotations, callouts) that have no project-management semantics
  \item Output formats prioritize fidelity (PDF, SVG) over convenience
\end{itemize}

\subsection{Deterministic Output}
\label{principle:determinism}

\textbf{Principle:} The same IR plus the same theme produces byte-identical output across all runs.

\textbf{Rationale:} Determinism enables:
\begin{itemize}
  \item Version control — changes to the IR produce predictable visual diffs
  \item CI/CD integration — automated rendering in pipelines
  \item Agent reliability — agents can trust their output
  \item Reproducibility — artifacts can be regenerated exactly
\end{itemize}

\textbf{Implications:}
\begin{itemize}
  \item No random layout algorithms; all positioning is computed deterministically
  \item Timestamps, UUIDs, and other entropy sources are excluded from output
  \item Font metrics must be stable (system fonts are problematic; embedded or specified fonts are required)
  \item Floating-point operations must be reproducible (or avoided)
\end{itemize}

\subsection{Human-Readable Source}
\label{principle:human-readable}

\textbf{Principle:} The IR must be readable and writable by humans without tooling.

\textbf{Rationale:} The IR serves dual audiences (humans and agents). Human readability ensures:
\begin{itemize}
  \item Onboarding friction is low — a PM can write a roadmap without learning a complex DSL
  \item Debugging is possible — when output is wrong, the IR is inspectable
  \item Manual overrides are feasible — humans can adjust agent-generated IR
  \item Version control diffs are meaningful
\end{itemize}

\textbf{Implications:}
\begin{itemize}
  \item YAML or YAML-like syntax preferred over JSON (less visual noise)
  \item Field names are descriptive English, not abbreviations
  \item Nesting depth is minimized
  \item Dates use ISO 8601; durations use human-readable formats where possible
\end{itemize}

\subsection{Agent-Friendly Authoring}
\label{principle:agent-friendly}

\textbf{Principle:} AI agents must be able to generate valid IR without special prompting or examples.

\textbf{Rationale:} A primary use case is agent-generated timelines. If agents struggle to produce valid IR, the system fails its purpose.

\textbf{Implications:}
\begin{itemize}
  \item The IR has a formal schema (JSON Schema or equivalent) that agents can reference
  \item Field names and structure align with how LLMs naturally express plans
  \item The IR avoids footguns — common mistakes should be invalid or auto-corrected
  \item Validation provides clear, actionable error messages
  \item The schema is small enough to fit in a context window
\end{itemize}

\subsection{Theme-Driven Rendering}
\label{principle:theme-driven}

\textbf{Principle:} Visual styling is entirely separated from content; all presentation decisions flow from the theme.

\textbf{Rationale:} Separation of content and presentation enables:
\begin{itemize}
  \item Rebranding without changing content — swap themes, regenerate
  \item Consistent organizational aesthetics across timelines
  \item Agents authoring content without making design decisions
  \item Themes as a distribution mechanism (theme marketplaces, branded themes)
\end{itemize}

\textbf{Implications:}
\begin{itemize}
  \item The IR contains no colors, fonts, or sizes — only semantic hints
  \item Themes define all visual properties
  \item The rendering model maps semantic IR elements to themed visual primitives
  \item Inline style overrides, if permitted, are explicit escapes from theme control
\end{itemize}

\subsection{Separation of Planning and Rendering}
\label{principle:separation}

\textbf{Principle:} The IR represents \emph{what to communicate}, not \emph{how to compute it}.

\textbf{Rationale:} Timeline Compiler is a rendering compiler, not a scheduling engine. The IR is a communication artifact, not a project plan. This separation:
\begin{itemize}
  \item Keeps the IR minimal and focused
  \item Avoids reimplementing project-management semantics
  \item Allows integration with any upstream planning tool
  \item Clarifies scope boundaries
\end{itemize}

\textbf{Implications:}
\begin{itemize}
  \item No dependency resolution, critical path calculation, or resource leveling
  \item Dates are provided, not computed
  \item Progress is asserted, not measured
  \item The IR describes the \emph{output} of planning, not the \emph{inputs} to planning
\end{itemize}

\subsection{Minimal Visual Clutter}
\label{principle:minimal-clutter}

\textbf{Principle:} Default output favors clarity over information density.

\textbf{Rationale:} Presentation-quality visuals communicate one message clearly rather than many messages densely. Executive audiences prefer clean visuals over comprehensive ones.

\textbf{Implications:}
\begin{itemize}
  \item Default themes use generous whitespace
  \item Labels are concise; detail lives in supporting materials
  \item Gridlines, axes, and chrome are minimized by default
  \item Information hierarchy is clear — not all elements receive equal visual weight
\end{itemize}

\subsection{Source-Control Friendly}
\label{principle:vcs-friendly}

\textbf{Principle:} The IR is designed for version control workflows.

\textbf{Rationale:} Modern teams version everything. If the IR produces noisy diffs or merge conflicts, adoption suffers.

\textbf{Implications:}
\begin{itemize}
  \item Text-based format (YAML preferred)
  \item Stable field ordering (or sorted keys)
  \item No generated IDs unless explicitly authored
  \item Changes to one element do not cascade spuriously to others
  \item Line-oriented structure where possible (each item on its own line)
\end{itemize}

\subsection{Composability}
\label{principle:composability}

\textbf{Principle:} Timeline definitions can reference, include, and compose other definitions.

\textbf{Rationale:} Real-world use involves:
\begin{itemize}
  \item Shared track definitions across timelines
  \item Theme inheritance and overrides
  \item Milestone libraries
  \item Multi-file organization for large timelines
\end{itemize}

\textbf{Implications:}
\begin{itemize}
  \item Support for file includes or references
  \item Namespacing to avoid collisions
  \item Clear resolution rules for overrides
  \item Themes can extend other themes
\end{itemize}

\subsection{Graceful Degradation}
\label{principle:graceful-degradation}

\textbf{Principle:} Invalid or incomplete IR produces useful output or clear errors, not silent failures.

\textbf{Rationale:} Authors (especially agents) make mistakes. The system should:
\begin{itemize}
  \item Validate early and report clearly
  \item Render partial output when possible (with warnings)
  \item Never produce corrupt or misleading visuals silently
\end{itemize}

\textbf{Implications:}
\begin{itemize}
  \item Schema validation runs before rendering
  \item Semantic validation catches contradictions (e.g., end before start)
  \item Warnings are visible but do not block output when the issue is cosmetic
  \item Errors specify exactly what is wrong and where
\end{itemize}

\subsection{Extensibility Without Complexity}
\label{principle:extensibility}

\textbf{Principle:} The IR and theme systems support extension without requiring core changes.

\textbf{Rationale:} Long-term success requires community contribution — new themes, new semantic elements, integration with new upstream tools. But extension points must not compromise core simplicity.

\textbf{Implications:}
\begin{itemize}
  \item Unknown IR fields are ignored (forward compatibility)
  \item Themes support custom properties via namespacing
  \item Plugin or extension architecture for renderers (future)
  \item Versioned schemas with clear migration paths
\end{itemize}

\subsection{Principle Summary}

\begin{table}[h]
\centering
\caption{Design Principles Summary}
\begin{tabularx}{\textwidth}{lX}
\toprule
\textbf{Principle} & \textbf{Core Commitment} \\
\midrule
Presentation-First & Optimize for board-ready visuals \\
Deterministic Output & Same input, same output, always \\
Human-Readable Source & Writable without tooling \\
Agent-Friendly Authoring & Agents generate valid IR naturally \\
Theme-Driven Rendering & Content and style fully separated \\
Separation of Planning and Rendering & Describe output, not compute schedule \\
Minimal Visual Clutter & Clarity over density \\
Source-Control Friendly & Clean diffs, easy merges \\
Composability & Reference, include, extend \\
Graceful Degradation & Clear errors, useful partial output \\
Extensibility Without Complexity & Extend without core changes \\
\bottomrule
\end{tabularx}
\end{table}
